\chapter{Dataset}

This study focuses exclusively on firms listed in United States equity markets. 
The U.S. market was selected for three principal reasons. 
First, it is the largest and most liquid equity market globally, providing a rich cross-section of firms across sectors and industries. 
Second, U.S. disclosure requirements are among the most stringent and standardized worldwide, ensuring high-quality, comparable corporate reporting. 
Third, the availability of structured financial, textual, and market data is substantially greater than in most international markets, making it particularly well-suited for large-scale empirical and machine learning research.

\section{SEC Filings}

Corporate disclosure in the United States is regulated by the U.S. Securities and Exchange Commission (SEC), which requires publicly traded firms to submit periodic filings through the Electronic Data Gathering, Analysis, and Retrieval (EDGAR) system. 
These filings serve as primary sources of information for investors, analysts, and regulators.

Among the wide range of SEC forms, this study focuses exclusively on Form 10-K and Form 10-Q. 
The 10-K is the firm's annual report, providing a comprehensive overview of financial performance, business operations, risk factors, and management analysis for the fiscal year. 
The 10-Q is the corresponding quarterly report, offering interim updates on financial condition and results of operations.
Together, these two forms represent the most detailed and standardized recurring disclosures made by U.S. public companies.

The choice to restrict attention to 10-K and 10-Q filings is motivated by their informational richness and regulatory consistency. 
Unlike event-driven filings (e.g., 8-K) that may be triggered by heterogeneous corporate events, 10-K and 10-Q reports are mandatory, regularly scheduled disclosures that follow standardized content requirements. 
Prior research demonstrates that these filings contain economically meaningful information relevant to valuation and risk assessment\cite{textual-analysis-dictionaries-and-10ks}. 
Focusing on these forms ensures comparability across firms and time, while concentrating on disclosures that are central to investors' expectations formation.

The textual data were obtained from the University of Notre Dame's cleaned SEC filings dataset\cite{sec_edgar_nd}, which removes HTML tags, tables, and formatting artifacts from raw EDGAR documents. 
This preprocessing is particularly important for natural language processing applications, as raw filings often contain extensive markup and tabular structures that are not directly usable in textual modeling.

The sample period spans 2011 to 2024. 
The starting point was selected to reflect modern market conditions and reporting practices. 
Following regulatory changes such as the adoption of XBRL (eXtensible Business Reporting Language) requirements in the late 2000s and early 2010s, SEC filings became more standardized and machine-readable. 
The SEC mandated phased XBRL adoption beginning in 2009, with full compliance across firms in the early 2010s, substantially increasing structural consistency in filings. 
Restricting the sample to 2011 onward therefore improves comparability and alignment with contemporary disclosure formats. 
Additionally, older financial data may reflect structural market conditions, such as lower algorithmic trading penetration and different regulatory environments, that are less representative of today's market microstructure.

\subsection{Section Extraction and Text Processing}

Although 10-K and 10-Q filings are information-rich, they are also lengthy and heterogeneous.
Not all sections are equally relevant for short-term stock price prediction. 
Therefore, rather than using the entire filing, this study focuses on two specific sections: 
Management's Discussion and Analysis of Financial Condition and Results of Operations (MD\&A) and Quantitative and Qualitative Disclosures About Market Risk.

The MD\&A section provides management's narrative explanation of financial performance, liquidity, capital resources, and operational trends. 
It frequently includes forward-looking statements, explanations of earnings changes, and contextual discussion that goes beyond raw financial statements. 
Prior literature highlights the importance of textual tone and narrative disclosures in influencing investor reactions.

The “Quantitative and Qualitative Disclosures About Market Risk” section addresses exposures to interest rate risk, currency risk, commodity price risk, and other financial uncertainties. 
These disclosures are directly relevant to risk assessment and may influence investors' expectations about future volatility and firm performance.

Section extraction was performed using regular expression matching to isolate these components from each filing. 
Because automated extraction may occasionally fail or capture unintended content, additional filtering steps were applied. 
Very short lines and lines without alphabetic characters were removed, as these often correspond to residual table fragments or formatting artifacts rather than meaningful narrative text. 
Samples with total character counts exceeding the 99th percentile were excluded, as these typically reflected extraction errors. 
Filings with fewer than two lines of extracted text, or with both targeted sections empty, were also removed to ensure substantive textual content (see Appendix \ref{appendix:quality_inspection} for quality inspection details
and Appendix \ref{appendix:text_length_statistics} for statistics).

\subsection{Firm Identifier Mapping: CIK to Ticker}

The cleaned SEC dataset contains firm identifiers in the form of Central Index Key (CIK) codes assigned by the SEC. 
However, stock price databases typically use ticker symbols. 
To align filings with market data, it was necessary to construct a mapping from CIK to ticker.

This mapping was created by combining two sources. 
First, the SEC's Insider Transactions Data Sets (Forms 3, 4, and 5) were used, as these filings often include both CIK and ticker information. 
These historical submissions provide time-varying mappings that are particularly useful for capturing changes in ticker symbols over time. 
Second, the SEC's publicly available ticker-to-CIK mapping file\cite{sec_ticker_list} was used to supplement recent mappings. 
Because the latter primarily reflects current ticker assignments and may not accurately represent historical changes, combining both sources improves coverage and historical accuracy.

After mapping CIK identifiers to stock tickers, each ticker is treated as a distinct firm in the dataset. This transformation allows us to move from a filing-level perspective to a firm-level panel structure analysis.

\subsection*{Overall Firm Coverage}

The final dataset contains 119,613 filings corresponding to 3,671 unique companies (see Appendix \ref{appendix:filings_per_company}). On average, each firm contributes 32.58 filings. This indicates that the dataset has both substantial cross-sectional breadth and meaningful longitudinal depth.

The relatively large number of firms ensures that the dataset captures heterogeneous disclosure behavior across industries and firm sizes. At the same time, the average of roughly 33 filings per firm suggests that many firms are observed repeatedly over time, making the dataset suitable for panel-based empirical analysis and machine learning models that benefit from repeated firm observations.

\section{Stock Price Data}

Stock price data were obtained from Stooq\cite{stooq_data}, which provides historical daily price data for companies listed on the NYSE and NASDAQ. 
The dataset includes daily open, high, low, close, and volume information.

To align financial reports with market data, each filing was treated as an event date. 
Rather than relying solely on price information from a single day, rolling windows of preceding trading days (e.g., 5, 10, 20 days) were used to compute summary statistics such as moving averages, momentum indicators, and realized volatility measures. 
This approach compresses multiple daily observations into a single feature vector associated with each filing event. 
By incorporating recent market dynamics prior to disclosure, the dataset captures prevailing trends and risk conditions that may interact with the informational content of the filing.

\subsection{Stock Price Feature Engineering}

For each filing event occurring at trading day $t$, a lookback window of up to $L=60$ trading days prior to $t$ is used. All features are computed exclusively on observations $\{t-L, \dots, t-1\}$.

The objective of this feature construction is threefold. First, it captures prevailing market trends prior to the disclosure. Second, it summarizes recent volatility and trading activity, which may interact with the informational content of the filing. Third, it provides economically interpretable baseline indicators against which the incremental predictive value of textual sentiment can be evaluated.

\subsubsection*{Notation}

Let $P_\tau$ denote the closing price at day $\tau$, $V_\tau$ the trading volume, $H_\tau$ the daily high price, and $L_\tau$ the daily low price. Define the simple daily return as:

\begin{equation}
r_\tau = \frac{P_\tau - P_{\tau-1}}{P_{\tau-1}},
\end{equation}

for $\tau \in \{t-L+1, \dots, t-1\}$.

All statistical quantities below are computed using only the historical sample $\mathcal{H}_t = \{t-L, \dots, t-1\}$.

\subsubsection*{Last Observed Price}

The most recent closing price prior to the filing is included as:

\begin{equation}
\text{Close}_{t-1} = P_{t-1}.
\end{equation}

Although returns are scale-invariant, the price level may still contain relevant information about firm characteristics, liquidity regimes, or price clustering effects.

\subsubsection*{Moving Averages}

To capture trend dynamics, simple moving averages (SMA) are computed over multiple horizons:

\begin{equation}
\text{MA}_k(t) = \frac{1}{k} \sum_{j=1}^{k} P_{t-j},
\end{equation}

for $k \in \{5, 10, 20, 50\}$, provided sufficient historical observations exist. If fewer than $k$ observations are available, an expanding mean over the available window is used:

\begin{equation}
\text{MA}_k(t) = \frac{1}{|\mathcal{H}_t|} \sum_{\tau \in \mathcal{H}_t} P_\tau.
\end{equation}

Moving averages are widely used in technical analysis to identify short and medium-term trends. Comparing the most recent price to these averages implicitly captures momentum and mean-reversion effects.

\subsubsection*{Momentum Indicators}

Momentum measures the cumulative return over a specified lag:

\begin{equation}
\text{Momentum}_k(t) = \frac{P_{t-1}}{P_{t-1-k}} - 1,
\end{equation}

for $k \in \{5, 20\}$. If fewer than $k$ lagged observations are available, the longest feasible lag within the historical window is used.

Momentum is included because empirical finance literature documents return continuation effects over short to intermediate horizons. Pre-filing momentum may influence how new information is interpreted or priced by the market.

\subsubsection*{Volatility Measures}

Realized volatility is approximated using the sample standard deviation of daily returns. For a window of length $k$, volatility is defined as:

\begin{equation}
\sigma_k(t) = 
\sqrt{
\frac{1}{k-1}
\sum_{j=1}^{k}
\left(r_{t-j} - \bar{r}_k\right)^2
},
\end{equation}

where $\bar{r}_k$ is the mean return over the same window. Volatility is computed for $k \in \{5, 20\}$, as well as over the full available lookback window:

\begin{equation}
\sigma_{\text{full}}(t) =
\sqrt{
\frac{1}{|\mathcal{H}_t|-1}
\sum_{\tau \in \mathcal{H}_t}
\left(r_\tau - \bar{r}\right)^2
}.
\end{equation}

Short-term volatility captures immediate market uncertainty prior to the filing, while longer-window volatility reflects the broader risk regime. Elevated pre-announcement volatility may signal heightened uncertainty or information asymmetry.

\subsubsection*{Volume Standardization}

To quantify abnormal trading activity, the most recent volume is standardized relative to its historical distribution:

\begin{equation}
Z^{\text{vol}}_t =
\frac{V_{t-1} - \mu_V}{\sigma_V},
\end{equation}

where $\mu_V$ and $\sigma_V$ denote the sample mean and standard deviation of volume over $\mathcal{H}_t$. This z-score measures whether trading activity immediately prior to the filing is unusually high or low. Abnormal volume may reflect information leakage, anticipation effects, or increased investor attention.

\subsubsection*{Intraday Price Range}

Finally, the high–low range of the last historical trading day is included as a proxy for intraday volatility:

\begin{equation}
\text{HLRange}_t =
\frac{H_{t-1} - L_{t-1}}{P_{t-1}}.
\end{equation}

This normalized range captures the intensity of price fluctuations within a single day, complementing return-based volatility measures.

\subsubsection*{Economic Interpretation}

Together, these features summarize multiple dimensions of pre-filing market conditions: trend (moving averages, momentum), risk (volatility measures, high-low range), and trading activity (volume z-score). Importantly, all quantities are computed strictly before the filing date, ensuring that the predictive models operate in a realistic, forward-looking setting.

Including these structured market features serves two purposes. First, they provide a benchmark level of predictive information derived purely from price dynamics. Second, they allow for an assessment of whether textual information extracted from financial disclosures adds incremental predictive power beyond what is already embedded in historical price and volume patterns.

\section{Sector and Industry Classification}

To enrich the dataset with structural firm characteristics, sector and industry classifications were incorporated using data from the NASDAQ Stock Screener\cite{nasdaq_screener}, which provides mappings from ticker symbols to sector and industry categories. 
Including sector and industry identifiers enables cross-sectional analysis of heterogeneous effects across economic segments and allows for later examination of sector-adjusted or industry-adjusted return behavior.

The dataset spans 12 unique sectors and 149 unique industries (see Appendix \ref{appendix:sector_industry_distribution}), indicating substantial economic diversity. 
The largest sectors by number of companies are Health Care, Consumer Discretionary, and Finance. In contrast, Basic Materials and Telecommunications are relatively underrepresented. This imbalance reflects broader patterns in U.S. equity markets, where certain sectors contain a larger number of publicly listed firms.

\section{Label Construction}

Since the objective of this study is to analyze the short-term impact of financial reports on stock price dynamics, we define prediction horizons ranging from one to seven trading days following the filing date. Formally, for each filing event occurring at trading day $i$, we construct labels corresponding to horizons $h \in \{1,2,\dots,7\}$. The task is therefore to predict, based solely on information available at the filing date, the stock price movement at day $i+h$.

Both regression and classification formulations are considered. In the regression setting, we focus on relative price changes rather than absolute price levels. This choice is motivated by the substantial heterogeneity in price magnitudes across firms. Predicting absolute price differences would introduce scale distortions, as a \$1 movement has different economic meaning for a \$10 stock compared to a \$500 stock. To ensure comparability across firms, labels are defined in terms of returns.

Specifically, for each horizon $h$, the regression label is computed as the simple return:

\begin{equation}
r_{i,h} = \frac{p_{i+h} - p_i}{p_i},
\end{equation}

where $p_i$ denotes the reference price at day $i$, and $p_{i+h}$ denotes the price at day $i+h$. This formulation captures the proportional price change over the specified horizon and provides a scale-invariant measure of market reaction.

In addition to regression, a classification formulation is implemented. Inspired by the approach of Lee et al.\cite{lee-etal-2014-importance}, we define a three-class labeling scheme: \textit{UP}, \textit{STAY}, and \textit{DOWN}. A sample is labeled as:

\begin{itemize}
    \item \textit{UP} if $r_{i,h} > 0.01$,
    \item \textit{DOWN} if $r_{i,h} < -0.01$,
    \item \textit{STAY} otherwise.
\end{itemize}

The 1\% threshold is chosen to filter out minor price fluctuations that may reflect normal market noise rather than economically meaningful reactions. In highly liquid U.S. equity markets, small daily movements often arise from microstructure effects, bid-ask bounce, or general volatility. By introducing a 1\% band around zero, the classification task focuses on substantial directional moves that are more likely to reflect genuine information incorporation following earnings disclosures.

\subsubsection*{Intraday Alignment Using Filing Acceptance Time}

To more precisely capture the immediate market reaction to disclosures, we incorporate the official SEC acceptance timestamp of each filing. Earnings reports may be released before market open, during trading hours, or after market close. Ignoring this timing dimension could misalign the return calculation and dilute the measured effect.

Accordingly, the reference price $p_i$ is defined conditional on the filing time:

\begin{itemize}
    \item If the filing is published \textit{before trading hours}, the return is computed between the previous day's closing price and the same day's opening price.
    \item If the filing is published \textit{during trading hours}, the return is computed between the opening and closing price of the same trading day.
    \item If the filing is published \textit{after trading hours}, the return is computed between the same day's closing price and the next trading day's opening price.
\end{itemize}

This alignment ensures that the computed return reflects the first tradable price adjustment following public disclosure. Consequently, the constructed labels better approximate the immediate information-driven component of price movement rather than unrelated market fluctuations.

\subsubsection*{Sector and Industry Adjusted Returns}

Beyond raw returns, additional sets of labels are constructed using sector and industry adjusted returns. The motivation for this adjustment is that stock price movements are often influenced not only by firm-specific information but also by broader sectoral or industry-wide trends. For instance, macroeconomic news affecting interest rates may systematically move financial stocks, while commodity price shocks may affect energy firms collectively. In such cases, part of the observed return may reflect common factors rather than firm-specific information contained in the filing.

To isolate the idiosyncratic component of the reaction, we compute adjusted returns by subtracting the contemporaneous mean return of firms within the same sector or industry. Formally, let $\mathcal{S}(j)$ denote the set of firms belonging to the same sector (or industry) as firm $j$. The adjusted return is defined as:

\begin{equation}
\tilde{r}_{i,h}^{(j)} = r_{i,h}^{(j)} - 
\frac{1}{|\mathcal{S}(j)| - 1}
\sum_{\substack{k \in \mathcal{S}(j) \\ k \neq j}} r_{i,h}^{(k)},
\end{equation}

where a leave-one-out approach is employed to avoid mechanical bias by excluding the firm's own return from the group average. This ensures that the benchmark represents the common movement of peer firms only.

The sector and industry adjusted labels allow for a more nuanced investigation of predictive performance. If a model performs well using raw returns but poorly using adjusted returns, this may indicate that it primarily captures sector-wide trends rather than firm-specific informational effects. Conversely, strong performance on adjusted labels would suggest that the model extracts idiosyncratic signals from financial disclosures that go beyond common industry movements.

By constructing both raw and adjusted labels across multiple horizons and prediction tasks, the dataset enables a comprehensive evaluation of short-term stock price reactions to financial reports under different definitions of market response.